% 原文 pp.47–52（印刷 pp.1165–1170）
\subsection{两种规范下的 BRST 上同调}
% 承接原文 §3.6（BRST 量子化的算符表述）：本段讨论 $A_0=0$ 与 $\Lambda=d^*A$ 两种规范，
% 随后进入 §3.7 反常。

% ===== 原文 p.47 =====
我们首先考虑 $u=0$ 的情形。此时拉格朗日量（Lagrangian）为
\[
\widehat L=\int d^4x\left(\frac{1}{4e^2}F^2+\frac{1}{2}v^2A_0^2-v\bar c\,\frac{Dc}{Dt}\right),
\]
其中 $D$ 表示协变导数（covariant derivative）。把 $\bar c$ 换成 $-v\bar c$，并让 $v$ 趋向无穷（这里用到「没有任何东西依赖于 $v$」这一事实），我们看到路径积分被局限到超平面 $A_0=0$ 上；在极限下便得到一个拉格朗日量
\begin{equation}
  \tag{3.8}
  \label{eq:3-8}
  \widehat L=\int d^4x\left(\frac{1}{4e^2}F^2+\bar c\,\frac{dc}{dt}\right).
\end{equation}

% [Tutor Note: 由第一式过渡到 (3.8) 的重标度必须算清楚。原式的鬼场项是 $-v\bar c\,Dc/Dt$，若按字面把 $\bar c$ 换成 $-v\bar c$，该项变成 $+v^2\bar c\,dc/dt$，系数为 $v^2$ 而非 $1$，与 (3.8) 不符；要得到 (3.8)，实际需要的替换是 $\bar c\to-\bar c/v$（其 Jacobi 因子 $\det(v^{-1})$ 只改变路径积分的整体归一化，而「没有量依赖 $v$」正说明该归一化可被吸收）。同一替换下 $\frac12v^2A_0^2$ 在 $v\to\infty$ 时给出 $\delta(A_0)$，这才是「路径积分被局限到 $A_0=0$」的确切含义——注意它是作为 $A_0$ 的泛函的 $\delta$ 函数，而不是把 $A_0$ 硬性设为 $0$。此处疑为原书排印之误，或作者略去了「同时把 $\widehat L$ 整体除以 $v^2$」这一步。]

既然 $A_0$ 现在为零，我们所得到的就是通常的量子力学，其中的动力学变量是一个空间联络、该联络的时间导数，以及鬼场（ghost）。因此 $\widehat{\mathcal H}$ 具有形式 $\mathcal O(\widetilde{\mathcal A})\otimes\Lambda\tilde{\mathfrak g}^{*}$，其中 $\widetilde{\mathcal A}$ 是空间闭路（spatial cycle）上的联络空间，$\mathcal O(\widetilde{\mathcal A})$ 是 $\widetilde{\mathcal A}$ 上的函数空间，$\tilde{\mathfrak g}$ 是空间闭路上规范变换群的 Lie 代数，而 $\Lambda\tilde{\mathfrak g}^{*}$（$\Lambda$ 表示外代数）是 $c\in\Pi\tilde{\mathfrak g}$ 的函数空间，这里 $\Pi$ 表示把 $\tilde{\mathfrak g}$ 当作纯粹奇的向量空间，即宇称反转。此外，鬼场的哈密顿量（Hamiltonian）为零，因为经典解空间上不存在非平凡的演化（鬼场的 Euler–Lagrange 方程不过是 $\frac{dc}{dt}=\frac{d\bar c}{dt}=0$）。于是我们理论中的哈密顿量就是通常的规范理论哈密顿量
\begin{equation}
  \tag{3.9}
  \label{eq:3-9}
  H=\int d^3x\left(\frac{1}{2e^2}F_A^2+\frac{e^2}{2}\nabla_A^2\right),\qquad A\in\widetilde{\mathcal A},
\end{equation}
它作用在张量积的第一个因子上。

% [Tutor Note: (3.9) 中 $\frac{1}{2e^2}F_A^2$ 只含空间分量，即磁场能 $\frac{1}{2e^2}B^2$；电场能已被 $\frac{e^2}{2}\nabla_A^2$ 取代。两个系数互为倒数并非笔误：正则动量是 $\pi_A=\frac{1}{e^2}\dot A$，故 $H=\int\left(\frac{e^2}{2}\pi_A^2+\frac{1}{2e^2}B^2\right)$，量子化 $\pi_A=-i\delta/\delta A$ 之后才出现 $\frac{e^2}{2}\nabla_A^2$。这里 $\nabla_A^2$ 按 $\delta^2/\delta A^2$ 型的（负半定）Laplace 算符理解，否则「电场能」的谱无下界。最常见的错误是把它写成 $\frac{1}{2e^2}\nabla_A^2$。]

容易看出，按鬼数（ghost number）分次的空间 $\widehat{\mathcal H}$ 不是别的，正是 Lie 代数 $\tilde{\mathfrak g}$ 以模 $\mathcal O(\widetilde{\mathcal A})$ 为系数的标准复形的空间。而且，由 $Q$ 的显式公式可以得知：在此情形下 $Q$ 恰好就是这个标准复形中的微分。于是，按定义为 $Q$ 的第 0 上同调（cohomology）的物理 Hilbert 空间 $H_{\rm phys}$，就是 $\tilde{\mathfrak g}$ 以 $\mathcal O(\widetilde{\mathcal A})$ 为系数的第 0 上同调。这不过是 $\mathcal O(\widetilde{\mathcal A})$ 中 $\tilde{\mathfrak g}$-不变量的空间，即 $\widetilde{\mathcal A}/G$ 上的函数空间，而后者按定义就是传统的 Hilbert 空间 $H_{\rm trad}$。

% [Tutor Note: 「$Q$ 恰是标准复形中的微分」背后的定理是：$\mathcal O(\widetilde{\mathcal A})\otimes\Lambda^{\bullet}\tilde{\mathfrak g}^{*}$ 正是 $\tilde{\mathfrak g}$ 以 $\mathcal O(\widetilde{\mathcal A})$ 为系数的 Chevalley–Eilenberg 复形 $C^{\bullet}(\tilde{\mathfrak g};\mathcal O(\widetilde{\mathcal A}))$，其微分为
% $(\delta_{CE}\omega)(X_0,\dots,X_n)=\sum_i(-1)^iX_i\,\omega(\dots,\widehat{X_i},\dots)+\sum_{i<j}(-1)^{i+j}\omega([X_i,X_j],\dots)$，
% 而 $X_i$ 通过 $A\mapsto A+\epsilon X_i$ 作用在 $\mathcal O(\widetilde{\mathcal A})$ 上。要验证 $Q=\delta_{CE}$，须用到 $\delta A=Dc$、$\delta c=-\frac12[c,c]$（即结构常数）以及 $\delta\bar c=\pi_{A_0}$ 在 $A_0=0$ 规范下退化为零这一点；换一个规范后 $\delta\bar c$ 不再为零，$Q$ 就会多出 (3.11) 中的 $Q'$。参见 Chevalley–Eilenberg 1948 年原文，或 Fuks《Cohomology of Infinite-Dimensional Lie Algebras》。]

这说明：BRST 上同调是 Lie 代数上同调的无穷维推广。

现在考虑由 $v=0,\ u=1$ 给出的另一种规范：$\Lambda=d^{*}A$。让我们用这个规范来看 $H_{\rm phys}$ 与 $H_{\rm trad}$ 之间的同构。我们有
\begin{equation}
  \tag{3.10}
  \label{eq:3-10}
  \widehat L=\int d^4x\left(\frac{1}{4e^2}F^2+\frac{1}{2}(d^{*}A)^2-\bar c\,D^{*}d_Ac\right).
\end{equation}

在此情形下，鬼场的运动方程是非平凡的，并且是二阶的；因此 Hilbert 空间 $\widehat{\mathcal H}$ 由 $A_s$、$c$、$\bar c$、$A_0$ 的函数构成，其中 $A_s$ 是空间闭路上的一个联络。此时人们发现
\begin{equation}
  \tag{3.11}
  \label{eq:3-11}
  Q=Q_{\text{Lie algebra cohomology}}+Q',
\end{equation}
其中 $Q'=\int\pi_{\bar c}\pi_{A_0}$，而 $\pi_{\bar c}$、$\pi_{A_0}$ 分别是 $\bar c$、$A_0$ 的动量算符。

直接验证 (3.11) 中两个被加项反对易、以及 $(Q')^2=0$，都是容易的。还可以验证：$Q'$ 除 0 次之外是无圈的（acyclic），在 0 次它有一维上同调。于是，由 Künneth 公式，我们再次得到 $H_{\rm phys}=H_{\rm trad}$。

% [Tutor Note: 此处把两种规范下的结果等同起来，用的是双复形的 Künneth 公式，其前提是 $(Q')^2=0$、$Q_{CE}^2=0$，且两者\textbf{反对易}（不是对易）；若误把 $\{Q_{CE},Q'\}$ 算成对易子，就会错误地得出 $Q^2\neq0$。$(Q')^2=0$ 的机制是 $\pi_{\bar c}$ 为奇量，故 $\pi_{\bar c}^2=0$（$\pi_{A_0}$ 为偶量，可以随意穿过）。所谓「$Q'$ 除 0 次外无圈」，即 $(A_0,\bar c)$ 构成同调代数意义下的一对可缩变量（contractible pair）：$Q'A_0=\pi_{\bar c}$、$Q'\bar c=0$，于是除常数外一切都被配对消掉。这套「平凡对」技巧在 Henneaux–Teitelboim《Quantization of Gauge Systems》中有系统处理。]

% ===== 原文 p.48 =====
\subsection{反常（Anomalies）}
% 原文 §3.7

现在让我们回顾 BRST 构造所必需的条件 1、2、3，并分析它们在何时得到满足。这些条件是

\begin{enumerate}
\item 拉格朗日量表述：$\delta^2=0$。\\
     哈密顿量表述：$Q^2=0$。
\item 拉格朗日量表述：$\delta$ 与 $\Lambda$、$\widetilde L$ 无关。\\
     哈密顿量表述：$Q$ 与 $\Lambda$、$\widetilde L$ 无关。
\item 拉格朗日量表述：$\delta\mathcal L_{\rm eff}=0$。\\
     哈密顿量表述：$[Q,H]=0$。
\end{enumerate}

% [Tutor Note: 这里的 $\widetilde L$ 与上一页（p.1165）的 $\widehat L$ 应指同一个对象，即经规范固定后、含鬼场项的拉格朗日量；原书在两页之间把帽子换成了波浪号，属记号不一致，读者不要误以为另有一个拉格朗日量。条件 2 的物理内容是：$\delta$（亦即 $Q$）不依赖于规范固定函数 $\Lambda$ 与 $\widetilde L$——这正是下面能把 (3.8) 与 (3.10) 两种规范下的 $H_{\rm phys}$ 直接比较、并断言上同调与规范选取无关的依据。]

正如我们已经提到的，性质 1 与性质 2 总是可以达到的。

然而，也正如我们已经提到的，若存在反常，性质 3 可能不成立。因此让我们更仔细地考察反常。

考虑一个四维规范场论，其中含一个手征旋量（chiral spinor）$\psi$，取值于规范群 $G$ 的一个复表示（complex representation）$\rho$，还含一个反手征旋量（antichiral spinor）$\bar\psi$，取值于对偶表示 $\bar\rho$。此类场的基本的规范不变拉格朗日量是
\begin{equation}
  \tag{3.12}
  \label{eq:3-12}
  \mathcal L=\frac{1}{4e^2}\int F^2+\int\bar\psi D_A\psi .
\end{equation}

在量子理论中我们关心的是路径积分 $\int e^{-\mathcal L}$。在这个路径积分中把 $\psi$ 积掉，我们得到
\begin{equation}
  \tag{3.13}
  \label{eq:3-13}
  \int\det(D_A)\,e^{-L(A)}\,DA,
\end{equation}
其中 $D_A:\Gamma(S_+)\to\Gamma(S_-)$ 是协变 Dirac 算符，而 $L(A)$ 是纯规范理论部分的拉格朗日量。

% [Tutor Note: (3.12) 中的 $\int F^2$ 应读作 $\int\mathrm{tr}(F\wedge*F)$（或 $\int\mathrm{tr}\,F_{\mu\nu}F^{\mu\nu}$），作者把迹与 Hodge 星号一并省略了；同理，欧氏号差下费米子项没有显式的 $i$，因为 $\bar\psi$ 与 $\psi$ 是彼此独立的 Grassmann 变量，而 $D_A$ 是反自伴的（见下文 (3.14) 后的括注）。若改用闵氏号差，则 $\bar\psi=\psi^{\dagger}\gamma^0$、该项带 $i$，(3.13) 中的 $\det D_A$ 也须相应换成 $\det(iD_A)$。]

积分 (3.13) 可能并不存在规范不变的调节（regularization）。更糟的是，它甚至可能不存在这样一种「非规范不变、但当截断趋于无穷时规范不变性得以恢复」的调节。在此情形下，我们所考虑的规范场论在量子力学意义上就讲不通，即使在微扰论中也讲不通，因为规范对称性无法在量子理论中得到尊重。这一现象称为反常（anomaly）。

反常有一个几何上的理由：虽然算符 $D_A$ 是规范不变的，但它的行列式 $\det(D_A)$ 一般不能定义为 $A$ 的规范不变函数。在通常的数学处理中，这个行列式不是联络的规范类空间上的函数，而是该空间上某个线丛的一个截面，这个线丛称为行列式线丛（determinant line bundle）；该丛带有一个典范联络（canonical connection）。如果这个典范联络不能把该丛平凡化，那么行列式就不是函数，而是一个非平凡丛的截面，因而不能有意义地积分。

把反常分成两类是有用的。

\begin{enumerate}
\item 局部反常（local anomaly）。典范联络具有非零曲率。此时，对适当选取的时空流形，这个曲率可以代表一个非平凡的第二上同调类，于是行列式线丛在拓扑上非平凡。
\item 整体反常（global anomaly）。典范联络是平坦的，但具有非平凡的单值性（monodromy）（并且丛也可能非平凡）。
\end{enumerate}

% ===== 原文 p.49 =====
因此，局部反常与整体反常都能产生拓扑性的反常，但只有前者能在微扰论中看到（办法是计算曲率）。

这里我们将只考虑局部反常。

\textbf{Remark 1.} 为了分析何时应当预期局部反常，可以从拓扑的观点来看待这个情形。我们假定时空 $M$ 是紧致的、可定向的（例如 $S^d$），并且指定了一个点 $\infty$，而所考虑的丛、联络与规范变换在无穷远处都平凡。在此情形下，联络的规范类（gauge classes）所构成的空间可以看作规范变换群 $\hat G$ 的分类空间（classifying space）$B\hat G$。$B\hat G$ 上的非平凡线丛由 $H^2(B\hat G)$ 分类。

现在，若 $M$ 紧致且可定向，我们就有嬗变映射（transgression map）$\tau:H_2(B\hat G)\to H_{d+2}(BG)$，其定义如下：给定一个二维同调类，取 $B\hat G$ 中代表它的一个曲面 $S$，并取相应的、$S$ 上的主 $\hat G$-丛。它的转移函数可以被看成 $(d+2)$ 维流形 $S\times M$ 上一个 $G$-丛的转移函数，这就定义了 $H_{d+2}(BG)$ 中的一个元素 $\tau_*([S])$。再考虑对偶映射 $\tau^{*}:H^{d+2}(BG)\to H^2(B\hat G)$。可以证明，我们这条线丛的 Chern 类就是 $\tau^{*}(C)$，其中 $C$ 是一个不依赖于 $M$ 的、固定的 $d+2$ 维特征类（characteristic class），它由曲率局部地算出。这个类恰好就是局部反常。

% [Tutor Note: 这段里 $G$ 与 $\hat G$ 必须分清：$G$ 是结构群，$\hat G$ 是「在无穷远处为恒等的」规范变换群。$\mathcal A/\mathcal G$ 同伦等价于 $B\hat G$，故行列式线丛的拓扑落在 $B\hat G$ 上；而反常作为特征类经嬗变映射 $\tau^{*}$ 来自 $BG$ 上的类 $C$——两个分类空间相差一次「时空方向的平移」（同调次数 $+d$），这正是 $H^2(B\hat G)$ 与 $H^{d+2}(BG)$ 的对应。$(S^3\mathfrak g)^{\mathfrak g}=H^6(BG)$ 是 Borel 定理的特例：对紧连通 Lie 群，$H^{\bullet}(BG;\mathbb R)\cong(S^{\bullet}\mathfrak g^{*})^G$，且 $S^k$ 对应 $H^{2k}$。]

于是，当 $d=4$ 时，局部反常位于 $H^6(BG)$ 中，或者位于 $(S^3\mathfrak g)^{\mathfrak g}$ 中，其中 $\mathfrak g$ 是 $G$ 的 Lie 代数。

例如，在粒子物理的标准模型中，规范群是 $SU(3)\times SU(2)\times U(1)$，因此反常的空间 $(S^3\mathfrak g)^{\mathfrak g}$ 是四维的：它等于四个子空间 $S^3_{\rm inv}(\mathfrak{su}(3))$、$S^3_{\rm inv}(\mathfrak{u}(1))$、$S^2_{\rm inv}(\mathfrak{su}(2))\otimes\mathfrak{u}(1)$、$S^2_{\rm inv}(\mathfrak{su}(3))\otimes\mathfrak{u}(1)$ 之和，而这四个子空间都是一维的（这里「inv」表示我们所取的是不变的对称多项式）。

这一讨论说明了：当所有费米子都处于规范群的实表示中时，为什么不会出现反常。事实上，在此情形下行列式线丛是实的，因而它的 Chern 类必定为零。

\textbf{Remark 2.} 尽管局部反常可以从上述拓扑观点来看待，人们仍应记住：它具有纯粹局部的本性，与时空的宏观结构无关。如果存在局部反常，那么量子理论在任何时空上都讲不通，即使在 $\mathbb{R}^d$ 上也讲不通。问题在于：即使行列式线丛在拓扑上平凡，它也不会有一个可以局部地定义出来的平坦联络；否则这个平坦联络对任何单连通的时空都适用，也就不会有拓扑反常了。因此，路径积分 (3.13) 即使在 $\mathbb{R}^d$ 上也是没有意义的。

\textbf{Remark 3.} 正如我们在上面 Remark 1 中所指出的，标准模型的规范群允许出现反常。然而，人们可以验证：来自标准模型各个不同物质场的反常会奇迹般地相互抵消。事实上，规范群在标准模型费米子上的表示（在添上若干无关紧要的实表示直和后）可以扩充为 $\mathrm{Spin}_{10}$ 的一个旋量表示；而对这个群，$H^6(BG)$ 为零，因此不可能有反常。在「大统一理论」中，人们正是这样来理解标准模型的：它由某个更统一的规范群（例如 $\mathrm{Spin}_{10}$）经规范对称性破缺而产生。

% [Tutor Note: 「$H^6(B\mathrm{Spin}_{10})$ 为零」的理由可以直接数生成元：$H^{\bullet}(B\mathrm{Spin}_{10};\mathbb R)$ 的低次生成元是 $p_1$（次数 4）、$p_2$（次数 8）以及 Euler 类 $e$（次数 10），次数 6 上没有任何东西，故 $H^6=0$。注意这里说的是「不可能有反常」，其含义是在这种旋量表示下三角反常系数恒为零；这与「标准模型内部各场反常相消」是同一件事的两种说法。]

现在让我们说明如何在微扰论中分析反常。我们的目标是解释：为什么在可能扩充场空间之后，性质 1 与

% ===== 原文 p.50 =====
性质 2（二者均见 §3.3）总是可以假定成立（即 $\delta^2=0$，且 $\delta$ 的定义不依赖于对某个特定拉格朗日量的选取），但人们不能假定有效拉格朗日量是 $\delta$-不变的。
% 原文在此跨页断句：「properties 1. and | 2. of section 3.3」，故 §3.3 一语在此补出。

首先，让我们仅仅试着在微扰论意义下理解积分 (3.13)。当我们写下 Feynman 图时，会发现单圈阶的发散无法以规范不变的方式消除。为了固定单圈阶，我们将通过添加另一个非常重的物质场 $\chi$ 来调节路径积分，使得它的行列式线丛与 $\psi$ 的行列式线丛互为逆。在有利的情形下，当 $\chi$ 的质量 $m$ 趋于无穷时，我们的原理论应当从这一理论中恢复出来。在另一些情形下，这一手续恰好会显示出反常为何存在。

为满足这一条件，物质场 $\chi$ 可以取为一个玻色型的场 $(\chi_+,\chi_-)$，取值于 $(S_+\oplus S_-)\otimes\rho$。此时复共轭场 $\bar\chi$ 取值于 $(S_+\oplus S_-)\otimes\bar\rho$，其中 $S_+$、$S_-$ 是 Poincaré 群的旋量表示（回想一下，在欧氏号差下 $S_+$ 与 $S_-$ 都既是自对偶的、又是自共轭的，即 $\bar S_{\pm}\cong S_{\pm}$）。不用说，这些场破坏了自旋–统计关系，因而和鬼场一样，不具有物理意义。

对场 $\chi_{\pm}$ 而言，自然的拉格朗日量项应当是
\begin{equation}
  \tag{3.14}
  \label{eq:3-14}
  L'(A,\chi_{\pm})=\int d^4x\Big((\bar\chi_+,D_A\chi_-)+(\bar\chi_-,D_A\chi_+)+m(\bar\chi_+,\chi_+)+m(\bar\chi_-,\chi_-)\Big).
\end{equation}
（这里 Dirac 算符是反自伴的。）

\textbf{Remark.} 这些 $\chi$ 场称为 Pauli–Villars 调节场（Pauli–Villars regulator fields）。

如果把表达式 (3.14) 加到拉格朗日量上，我们得到的将是行列式的模平方，而不是行列式本身，因此并不能修正反常。于是我们以一种破坏规范不变性的方式来修改 (3.14)：令 $A_0$ 为一个固定的联络，并置
\begin{equation}
  \tag{3.15}
  \label{eq:3-15}
  L''(A,\chi_{\pm})=\int d^4x\Big((\bar\chi_+,D_A\chi_-)+(\bar\chi_-,D_{A_0}\chi_+)+m(\bar\chi_+,\chi_+)+m(\bar\chi_-,\chi_-)\Big).
\end{equation}

现在考虑拉格朗日量为 $\mathcal L+L''$ 的理论。把 $\chi$ 场积掉，我们会得到一个因子 $\det(D_AD_{A_0}-m^2)^{-1}$。当 $m=0$ 时，这个因子在相差一个乘性常数的意义下是规范不变的，并且抵消分子中由反常引起的那个行列式。这说明：在这个理论中，对任何有限的 $m$ 都不存在拓扑反常（也就是说，相应的行列式线丛是平凡的）。然而，当 $m>0$ 时，规范不变性失效了。因此我们必须研究 $m\to\infty$ 的极限（按设想它应当给出我们的原理论），看看规范对称性是否会重新出现。

把行列式之比 $\det(D_AD_{A_0}-m^2)/\det(D_A)$ 沿规范变换 $t\in\hat{\mathfrak g}$ 的方向求导，我们得到（利用路径积分的解释）：它等于 $m\big\langle\int[(\bar\chi_+,t\chi_+)+(\bar\chi_-,t\chi_-)]\big\rangle$，其中 $\langle X\rangle$ 表示 $X$ 的期望值。这个期望值可以按 $1/m$ 的幂次展开。

为看出规范不变性的失效在 $m\to\infty$ 时是否持续存在，让我们考虑曲率算符 $F$ 的两点函数。容易看出，这个函数沿 $t$ 方向求导后的领头贡献（就 $1/m$ 而言）来自这样一张单圈图：一个 $\chi$ 圈，其内部插入 $t$ 算符，并有两条外出的 $F$-边。这一贡献是 $1/m$ 的零阶项，其形式为 $\sum d_{abc}t^aF^bF^c$，其中 $d_{abc}$ 是某个张量。因此，如果 $d_{abc}\neq0$，则规范不变性的失效在极限下依然残留。

% [Tutor Note: 「$\chi$ 圈内插入 $t$、两条外线为 $F$」正是三角图。此处最易被忽略的一步是量纲：质量项 $m(\bar\chi,\chi)$ 破坏规范不变性的效应形式上被 $1/m$ 压低，但三角图是线性发散的，发散度恰好补偿了这个 $1/m$，于是 $m\to\infty$ 时残留一个 $m^0$ 阶的 $\sum d_{abc}t^aF^bF^c$。也就是说，「重调节场应当跑掉、只留下可重整化的规范理论」这一直觉，只在 $d_{abc}=0$ 时才成立。这里 $d_{abc}$ 即 $\mathrm{tr}\big(T_a\{T_b,T_c\}\big)$ 型的对称不变量，其消失等价于该表示没有三角反常。]

% ===== 原文 p.51 =====
\textbf{Remark.} 如果原来的费米子处于实表示中，则 $d_{abc}$ 为零，(3.14) 中的调节是完全令人满意的。问题出现在原表示是复表示的时候。此时，除非放弃规范不变性，形如 (3.14) 的调节就不管用；即使 $d_{abc}=0$ 也是如此。标准模型就是一个例子：其中的费米子处于复表示，$d_{abc}=0$，而沿着上述思路并不存在规范不变的调节。

\textbf{Remark.} 当 $d_{abc}=0$ 时，不存在更高圈（或非微扰）的规范不变性破坏。

现在让我们从 BRST 的观点来考察反常。如果存在局部反常，我们就会有 $U=\delta\mathcal L_{\rm eff}(A)\neq0$（这里 $\mathcal L_{\rm eff}(A)$ 是把鬼场积掉之后、关于 $A$ 的有效拉格朗日量）。然而，由于反常是局部的，$U$ 必须是 $A$ 与 $c$ 的某个局部表达式的积分，并且该表达式关于 $c$ 是线性的。同样清楚的是 $\delta U=0$。

另一方面，如果 $U=\delta N$，其中 $N$ 是 $A$ 的某个局部表达式的积分，那么我们就可以通过把拉格朗日量重新定义为 $\mathcal L\to\mathcal L+N$ 来使 $\delta\mathcal L_{\rm eff}=0$。因此，我们应当定义 $\delta$ 的一个上同调：它作用在由规范场、物质场（若有）与鬼场构成的局部泛函上，并模去全微分（total derivatives）。反常位于 $\delta$ 的 1 次（也就是鬼数 1）上同调之中。

现在让我们说明：Lie 代数 $\mathfrak g$ 上的不变对称张量（等价地，$BG$ 的上同调类）可以被用来构造这个上同调中的 1-上闭链（1-cocycle）。反过来（这些类生成 1 次上同调）也是对的，但我们不去证明它。

设 $C$ 是 $S^{n+1}\mathfrak g$ 中的一个 $G$-不变元素。与这个元素相对应的是 $G$-丛的一个 $(2n+2)$ 维特征类，即 $C(F^{n+1})$，其中 $F$ 是曲率。与 $C$ 相对应的 Chern–Simons 形式 $CS_C(A)$ 是这样的局部 $(2n+1)$ 形式：
\[
\frac{\widetilde\delta\,C(F^{n+1})}{\widetilde\delta A}=CS_C(A)\wedge\widetilde\delta A,
\]
其中的等号准确到局部形式的全微分；这里 $\widetilde\delta$ 表示变分，以区别于 BRST 微分 $\delta$。

Chern–Simons 形式的主要性质如下：虽然这个形式不是规范不变的，但沿一个无穷小规范变换取它的 Lie 导数，得到一个局部形式的全微分。

% [Tutor Note: 这里的 $\widetilde\delta$ 是「对场 $A$ 的变分」，与 BRST 微分 $\delta$ 必须严格区分；原文用两种不同的符号正是为此，抄录时极易混同。上式就是通常的下降方程（descent equations）的起点：由 $d\,CS_C(A)=C(F^{n+1})$ 与 $\widetilde\delta CS_C=d(\cdots)$ 出发，把 $\widetilde\delta A$ 换成鬼场 $c$（即令 $\widetilde\delta\to\delta$、$\widetilde\delta A\to Dc$），就得到 $\delta CS_C=W(A,c)$。归一化要当心：$CS_C$ 中含有 $1/(n+1)$ 的因子（例如 $n=1$ 时 $CS=\mathrm{tr}\big(A\,dA+\tfrac23A^3\big)$），漏掉它会使 $W$ 的系数、从而反常系数相差 $n+1$ 倍。]

现在令 $M^{2n}$ 为我们的时空。设 $A$ 是 $M^{2n}$ 上的一个联络。我们想定义一个形如 $U(A,c)=\int W(A,c)$ 的泛函，其中 $W(A,c)$ 是 $M^{2n}$ 上的一个 $2n$ 形式，它对 $A$ 与 $c$ 是局部的，但不是规范不变的，并且满足 $\delta W=0$。$W$ 关于 $c$ 将是线性的。

设 $X^{2n+1}$ 是一个光滑流形，其边界等于 $M^{2n}$。以任意方式选取联络 $A$ 到 $X^{2n+1}$ 上的一个扩充（为简单起见，我们假定对 $X$ 的选取以及 $A$ 的扩充都不存在拓扑障碍；这个假定其实是无关紧要的）。然后置 $V(A)=\int_X CS_C(A)$。这个泛函依赖于 $A$ 到 $X$ 的扩充。然而，由 $CS$ 的主要性质，泛函 $\delta V(A)=W(A,c)$（其中 $\delta$ 是 BRST 微分）并不依赖于该扩充，因而它是 $A$ 与 $c$ 的一个局部泛函，且关于 $c$ 线性。于是 $U=\int W(A,c)$ 代表了局部泛函的 $\delta$ 上同调中的一个上同调类；可以证明它是非零的。这样我们就得到一个单射（事实上是同构）
\[
S^{n+1}(\mathfrak g)^{\mathfrak g}\longrightarrow H^{1,\,local}_{\delta}.
\]
对于四维理论，$n=2$，而上闭链来自 $(S^3\mathfrak g)^{\mathfrak g}=H^6(BG)$。

% [Tutor Note: 「可以证明它非零」以及后面的「事实上是同构」，其依据是 Gel'fand–Fuks 关于形式向量场 Lie 代数上同调的计算与 Stora 的下降方程方法：从不变多项式空间到鬼数 1 的局部上同调的映射是单射，而在 $d=4$ 时它还是满射。这条同构正是把「反常」等同于「$H^6(BG)$ 中的一个类」的严格依据；细节可参见 R. Stora 关于下降方程的讲义（1983），以及 Fuks《Cohomology of Infinite-Dimensional Lie Algebras》。]

% ===== 原文 p.52 =====
\subsection{参考文献与交叉引用}

关于规范场论 BRST 量子化的另一个文献来源，是 [II-Strings] 的第 4 讲与 [I-Faddeev] 的第 5 讲。经典规范场论见 [I-Classical Fields] 第 4 章。

% [Tutor Note: 方括号中的代号是本丛书（Quantum Fields and Strings: A Course for Mathematicians, AMS, 1999）内部的交叉引用：「II-Strings」指第 II 卷中关于弦论的部分，「I-Faddeev」指第 I 卷中 Faddeev 撰写的量子化章节，「I-Classical Fields」指第 I 卷中经典场论章节。译文一律原样保留这些代号，以便与原文对照。]

关于 BRST 量子化的入门材料，例如可参见

\begin{quote}
S. Weinberg, \textit{The Quantum Theory Of Fields}, chapter 15, vol. 2, Cambridge University Press, 1996.
\end{quote}

关于反常的入门材料也可以在同一文献中找到。
